\documentclass[10pt,letterpaper]{article}

\usepackage[T1]{fontenc}
\usepackage[utf8]{inputenc}
\usepackage[spanish,es-noquoting,es-noshorthands]{babel}
\usepackage{lmodern}
\usepackage[margin=1.9cm,top=1.5cm,bottom=1.6cm]{geometry}
\usepackage{booktabs}
\usepackage{array}
\usepackage{xcolor}
\usepackage{graphicx}
\usepackage{fancyvrb}
\usepackage{tikz}

\definecolor{acc}{HTML}{1F4E79}
\definecolor{gris}{HTML}{5B6672}

% Marca de algo pendiente (sale en rojo). Antes de entregar no debe quedar ningún uso.
\newcommand{\falta}[1]{{\color{red}\bfseries [FALTA: #1]}}

\setlength{\parskip}{2pt}
\setlength{\topsep}{2pt}
\setlength{\partopsep}{0pt}
\setlength{\parindent}{0pt}
\renewcommand{\arraystretch}{1.05}
\setlength{\emergencystretch}{4em}

% Secciones y pie de pagina sin paquetes externos: solo LaTeX base.
\makeatletter
\renewcommand\section{\@startsection{section}{1}{\z@}%
  {7pt \@plus 2pt \@minus 1pt}{3pt}%
  {\normalfont\normalsize\bfseries\color{acc}}}
\renewcommand\subsection{\@startsection{subsection}{2}{\z@}%
  {5pt \@plus 1pt}{2pt}%
  {\normalfont\small\bfseries\color{acc}}}
\renewcommand\@seccntformat[1]{\csname the#1\endcsname.\quad}
\renewcommand{\ps@plain}{%
  \renewcommand{\@oddhead}{}\renewcommand{\@evenhead}{}%
  \renewcommand{\@oddfoot}{\footnotesize\color{gris}%
    TI-4601 $\cdot$ Laboratorio 3 $\cdot$ Conjuntos de réplicas en MongoDB
    \hfill \thepage}%
  \renewcommand{\@evenfoot}{\@oddfoot}}
\makeatother
\pagestyle{plain}

\fvset{fontsize=\footnotesize,frame=leftline,framerule=0.6pt,rulecolor=\color{gris},xleftmargin=6pt}

\begin{document}

{\LARGE\bfseries\color{acc} Laboratorio 3 --- Conjuntos de réplicas en MongoDB}\\[2pt]
{\footnotesize\color{gris} TI-4601 Bases de Datos Avanzadas $\cdot$ Grupo 3 $\cdot$ Fabricio Mena $\cdot$ Christopher Rodriguez $\cdot$ Alisson Redondo}

\vspace{4pt}\hrule\vspace{6pt}

\section{Configuración del conjunto (pasos 0 a 3)}

\textbf{Máquina de los pasos 0 a 6.} Laptop AMD Ryzen 7 5700U, 16 GB de RAM (4 CPU y 9,7 GiB para WSL 2),
Windows 11 Home con Ubuntu 24.04 en WSL 2, Docker Desktop 29.5.2 y MongoDB 8.0.34. Registro con hora y
\texttt{OK} de cada paso en \texttt{evidence/comandos.md}.

\textbf{Paso 1: Compose.} Los tres nodos arrancan con el mismo comando:
\begin{Verbatim}
mongod --replSet rs-lab3 --bind_ip_all --wiredTigerCacheSizeGB 0.25
\end{Verbatim}

\begin{center}\footnotesize
\begin{tabular}{@{}lllll@{}}
\toprule
\textbf{Servicio} & \textbf{\texttt{-{}-replSet}} & \textbf{\texttt{hostname}} & \textbf{Volumen} & \textbf{¿Publica puertos?} \\
\midrule
\texttt{mongo-1} & \texttt{rs-lab3} & \texttt{mongo-1} & \texttt{mongo1} & no \\
\texttt{mongo-2} & \texttt{rs-lab3} & \texttt{mongo-2} & \texttt{mongo2} & no \\
\texttt{mongo-3} & \texttt{rs-lab3} & \texttt{mongo-3} & \texttt{mongo3} & no \\
\bottomrule
\end{tabular}
\end{center}

Tras \texttt{make lab3-up} los tres contenedores quedaron \texttt{healthy} y \texttt{make lab3-status}
respondió \texttt{NotYetInitialized} (\texttt{cluster-start.txt}).
\texttt{-{}-replSet} solo dice a qué conjunto puede pertenecer cada proceso; quiénes son los miembros y
quién vota se define con \texttt{rs.initiate()}, y hasta entonces no hay conjunto. Por eso el \texttt{\_id}
de la configuración tiene que ser \texttt{rs-lab3}: \texttt{mongod} rechaza una configuración con otro
nombre. Los miembros se nombran \texttt{mongo-1:27017} y no \texttt{localhost:27017} porque la
configuración se copia a los tres nodos y cada uno la usa para llegar a los otros: dentro de un contenedor
\texttt{localhost} es él mismo, y los nombres \texttt{mongo-N} los resuelve el DNS de la red de Compose.
Si se borran los volúmenes se pierden los datos y también la configuración del conjunto, que cada nodo
guarda en su base \texttt{local}; habría que repetir desde el paso 1.

\textbf{Paso 2: \texttt{rs.initiate()}.}

\begin{center}
\begin{tikzpicture}[nodo/.style={draw,rounded corners=2pt,minimum width=2.7cm,minimum height=0.85cm,align=center,font=\footnotesize},>=stealth]
\node[nodo,fill=acc!12] (m1) at (0,0) {\textbf{mongo-1:27017}\\PRIMARY};
\node[nodo] (m2) at (-3.2,-1.5) {\textbf{mongo-2:27017}\\SECONDARY};
\node[nodo] (m3) at (3.2,-1.5) {\textbf{mongo-3:27017}\\SECONDARY};
\node[nodo,fill=gris!10] (app) at (-6.2,0) {\textbf{app-mongo}\\pymongo};
\draw[->,thick] (m1) -- node[above left,font=\scriptsize]{oplog} (m2);
\draw[->,thick] (m1) -- node[above right,font=\scriptsize]{oplog} (m3);
\draw[<->,dashed,gris] (m2) -- node[below,font=\scriptsize]{latidos cada 2 s entre los tres} (m3);
\draw[->] (app) -- node[above,font=\scriptsize]{escrituras} (m1);
\end{tikzpicture}
\end{center}

\vspace{-2pt}
\begin{Verbatim}
rs.initiate({
  _id: "rs-lab3",
  members: [
    { _id: 0, host: "mongo-1:27017" },
    { _id: 1, host: "mongo-2:27017" },
    { _id: 2, host: "mongo-3:27017" }
  ]
})
\end{Verbatim}

\texttt{\_id: "rs-lab3"} es el nombre del conjunto y debe ser igual al \texttt{-{}-replSet}.
\texttt{members} es la lista de miembros: en cada uno, \texttt{\_id} es su número dentro del conjunto (0, 1
y 2) y \texttt{host} es el nombre y puerto con el que los demás lo alcanzan. Respondió \texttt{\{ ok: 1 \}}
y en unos segundos quedaron \texttt{mongo-1} como \texttt{PRIMARY} y \texttt{mongo-2} y \texttt{mongo-3}
como \texttt{SECONDARY}.

Se usa un número impar de votantes porque la mayoría es $\lfloor n/2 \rfloor + 1$: con 3 votantes es 2
y se tolera una caída; con 4 la mayoría es 3 y se sigue tolerando solo una, así que el cuarto nodo no
agrega tolerancia y en una partición de 2 contra 2 nadie tendría mayoría. MongoDB completó solo los campos
\texttt{arbiterOnly: false}, \texttt{buildIndexes: true}, \texttt{hidden: false}, \texttt{priority: 1},
\texttt{tags: \{\}}, \texttt{secondaryDelaySecs: 0} y \texttt{votes: 1} de cada miembro, y los
\texttt{settings} del conjunto (\texttt{rs-inspect.txt}). Con \texttt{localhost} cada nodo intentaría
conectarse a sí mismo en vez de a los otros, y el conjunto no se formaría.

\textbf{Paso 3: lectura de \texttt{rs.status()} y \texttt{rs.conf()}} (\texttt{rs-inspect.txt}).
El primario es \texttt{mongo-1}, y \texttt{mongo-2} y \texttt{mongo-3} copian de él
(\texttt{syncSourceHost: mongo-1:27017}). El \texttt{term} es 1, porque en un conjunto nuevo solo hubo
una elección. Los tres tienen \texttt{votes: 1} y \texttt{priority: 1}. En \texttt{settings},
\texttt{heartbeatIntervalMillis} es 2\,000 y \texttt{electionTimeoutMillis} es 10\,000: cada nodo manda un
latido cada 2 s, y si un secundario pasa 10 s sin saber del primario (unos cinco latidos perdidos) pide
votos. \texttt{db.hello()} en \texttt{mongo-1} da \texttt{setName: rs-lab3},
\texttt{isWritablePrimary: true} y \texttt{secondary: false}.

Entre que el primario cae y un secundario pide votos pasan como máximo unos 10 s: ese plazo es
\texttt{electionTimeoutMillis}. El \texttt{term} equivale al término de Raft del Lab 1: sube en cada
elección, y un primario que ve un \texttt{term} mayor deja de serlo. \texttt{priority} decide entre
candidatos que están al día: gana preferencia el de mayor prioridad, y uno con \texttt{priority: 0} nunca
será primario. Aquí los tres tienen 1, así que cualquiera puede ganar.

\section{Datos, carga e índices (pasos 4 y 5)}

\textbf{Predicción} (\texttt{evidence/predicciones.md}, escrita antes de correr): cada documento quedaría
en los 3 nodos, y 2 (el primario y un secundario) deberían tenerlo antes de que \texttt{insert\_many}
responda, porque la conexión usa \texttt{w: "majority"} y la mayoría de 3 es 2. Se cumplió:
\texttt{generar\_datos.py} produjo 400 eventos, 1\,189 zonas y 162\,810 de capacidad, con SHA-256
\texttt{9fc6a2db\ldots 774a} (\texttt{datos-manifest.txt}); \texttt{cargar.py} dio sus cinco líneas
\texttt{OK:} y \texttt{make lab3-check} terminó en \texttt{6/6}, con diferencia máxima de \emph{optime}
entre miembros de 0,0 s (\texttt{config-check.txt}).

\textbf{TODO 1.} \texttt{cargar\_eventos} lee el JSONL línea por línea, junta lotes de 500 y los inserta
con \texttt{insert\_many(..., ordered=False)}, conservando el \texttt{\_id} de cada evento. Con
\texttt{ordered=False}, si un documento falla (por ejemplo, un \texttt{\_id} repetido) el servidor sigue
con el resto del lote y reporta los errores al final; con \texttt{ordered=True} se detendría en el primero y
dejaría el lote a medias. Cada \texttt{insert\_many} espera a que 2 de los 3 nodos tengan el lote; con 400
eventos fue una sola llamada.

\textbf{TODO 2.} Los índices son \texttt{\{categoria: 1, fecha: 1\}} en \texttt{eventos} y
\texttt{\{evento\_id: 1, estado: 1\}} en \texttt{ordenes}: primero el campo de igualdad y después el de
rango u orden. Con \texttt{categoria} primero, los eventos de una categoría quedan juntos en el índice y
ya ordenados por \texttt{fecha}, así que el rango \texttt{fecha >= d} es un solo tramo y el
\texttt{sort} sale del índice sin ordenar en memoria.

Consulta: \texttt{categoria: "concierto", fecha >= "2026-12-01"}, ordenada por fecha.

\begin{center}\footnotesize
\begin{tabular}{@{}lll@{}}
\toprule
 & \textbf{Con índice} & \textbf{Sin índice} (\texttt{hint \$natural}) \\
\midrule
Etapas del plan & \texttt{IXSCAN} \texttt{categoria\_1\_fecha\_1} $\to$ \texttt{FETCH} & \texttt{COLLSCAN} $\to$ \texttt{SORT} \\
\texttt{totalKeysExamined} & 67 & 0 \\
\texttt{totalDocsExamined} & 67 & 400 \\
\texttt{nReturned} & 67 & 67 \\
\texttt{executionTimeMillis} & 1 & 1 \\
\bottomrule
\end{tabular}
\end{center}

Las dos dan los mismos 67 eventos (\texttt{explain-indices.txt}). Con el índice se examinaron solo esos
67 documentos y no hubo etapa \texttt{SORT}; sin índice se recorrieron los 400 y hubo que ordenar en
memoria. El tiempo empata en 1 ms porque 400 documentos son pocos, pero los documentos examinados sin
índice crecen con la colección y con el índice crecen solo con el resultado.

\section{Replicación observada (pasos 6 y 7)}

\textbf{Paso 6: oplog} (\texttt{oplog.txt}). La orden \texttt{ORD-MANUAL-1} insertada en el primario
aparece en \texttt{local.oplog.rs} con \texttt{op: 'i'} (insert), \texttt{ns: 'boleteria.ordenes'}
(base y colección), \texttt{o} (el documento insertado), \texttt{ts: Timestamp(\{ t: 1792700740, i: 1 \})}
(segundo y contador dentro de ese segundo, que ordenan las operaciones) y \texttt{t: 1} (el \texttt{term}
del primario que la escribió). Los dos secundarios estaban a 0 s del primario y \texttt{mongo-2} ya podía
leer la orden.

El oplog es una colección capada: tiene tamaño fijo (aquí 47\,287 MB configurados) y cuando se llena
sobrescribe las entradas más viejas, para que la replicación no crezca sin límite. La ventana
(\texttt{log length start to end}) era de 151 s porque el conjunto era nuevo. Si un secundario se atrasa
más que la ventana, las entradas que le faltan ya se borraron, no puede ponerse al día con el oplog y
necesita un \emph{initial sync}, que copia otra vez todos los datos.

Al escribir en \texttt{mongo-2} la respuesta fue \texttt{NotWritablePrimary (code 10107): not primary}:
solo el primario acepta escrituras y los secundarios solo aplican lo que trae el oplog. El cliente tiene
que encontrar al primario, y por eso la URI del driver lista los tres miembros.

\textbf{Paso 7: lecturas en secundario} (\texttt{lectura-secundario.csv},
\texttt{lectura-secundario.txt}). Este paso se ejecutó en una ASUS Vivobook X1504ZA (Intel Core i7-1255U, Windows con Docker Desktop y
WSL 2, kernel 6.18.54.1, MongoDB 8.0.34), con 500 intentos por modalidad.
\textbf{Predicción:} \falta{predicción del Paso 7, escrita antes de ejecutar (Persona B)}

\begin{center}\footnotesize
\begin{tabular}{@{}lrrr@{}}
\toprule
\textbf{Modalidad} & \textbf{Intentos} & \textbf{Obsoletas} & \textbf{Porcentaje} \\
\midrule
\texttt{w1\_secundario} & 500 & 499 & 99,8\,\% \\
\texttt{majority\_secundario} & 500 & 3 & 0,6\,\% \\
\texttt{causal\_majority} & 500 & 0 & 0,0\,\% \\
\bottomrule
\end{tabular}
\end{center}

Las lecturas obsoletas fueron 499/500 con \texttt{w:1}, pues el primario no espera
replicación; 3/500 con \texttt{majority}, ya que un tercer nodo puede seguir
atrasado aunque dos confirmen; y 0/500 con sesión causal y garantías
\texttt{majority}, que permite observar escrituras previas de esa sesión.
Esto no garantiza frescura para lecturas arbitrarias desde un secundario.

\textbf{CAP/PACELC.} Sin quórum, MongoDB no puede confirmar escrituras por
mayoría: privilegia su consistencia frente a la disponibilidad de esas
escrituras (CAP). Sin partición, \texttt{w:1} reduce la espera y
\texttt{majority} exige replicación (PACELC). Para ``mis entradas'' se
prefiere lectura causal o del primario; ``eventos destacados'' puede tolerar
secundarios atrasados. Los tres nodos comparten una computadora: estos
porcentajes no se extrapolan a una red distribuida real.


\section{Atomicidad y costo de las garantías (pasos 8 y 9)}

\textbf{Paso 8: compra atómica} (\texttt{comprar.txt} y
\texttt{comprar-roto.txt}). El \textbf{TODO 3} registra cada orden con
\texttt{\_id} único y estado \texttt{pendiente}; ante
\texttt{DuplicateKeyError} devuelve el estado ya guardado y no vuelve a
descontar. El stock se actualiza con un solo \texttt{update\_one}:
\texttt{\$elemMatch} exige \emph{en el mismo elemento} del arreglo
\texttt{zonas} el nombre pedido y \texttt{disponibles >= cantidad};
\texttt{\$inc} sobre \texttt{zonas.\$.disponibles} efectúa el descuento
atómico. Dos condiciones independientes, \texttt{zonas.nombre} y
\texttt{zonas.disponibles}, podrían satisfacerse en zonas distintas
(p.~ej., General con 0 y VIP con 10), admitiendo una compra inválida.

\begin{center}\footnotesize
\begin{tabular}{@{}lrr@{}}
\toprule
\textbf{Resultado (200 compradores; 120 entradas)} & \textbf{Correcta} & \textbf{Rota} \\
\midrule
Órdenes confirmadas / agotadas & 120 / 80 & 200 / 0 \\
Entradas General restantes & 0 & 115 \\
VIP: compras confirmadas de 4 sobre 10 & 2 & 3 \\
\bottomrule
\end{tabular}
\end{center}

La versión rota separó \texttt{find\_one} y \texttt{update\_one} con
\texttt{\$set}: varias solicitudes leyeron el mismo saldo y sobrescribieron
las actualizaciones de otras (\emph{lost updates}). En cambio, la versión
atómica pasó las seis verificaciones, incluida la idempotencia (cinco
reintentos de la misma orden, un solo descuento). La atomicidad cubre
\emph{un documento}, no toda la compra: si el proceso cae entre descontar
el evento y marcar la orden \texttt{confirmada}, puede quedar
\texttt{pendiente} con stock ya reducido. Se detectarían órdenes pendientes
antiguas para reconciliarlas; no se debe reejecutar ciegamente el descuento.

\textbf{Paso 9: costo de \emph{write concern}} (\texttt{write-concern.csv},
\texttt{write-concern.txt}). El \textbf{TODO 4} usa
\texttt{with\_options} y \texttt{WriteConcern} para configurar la garantía; mide cada
\texttt{insert\_one} con \texttt{perf\_counter}, sin incluir la creación
del cliente en la latencia. Hubo 30 inserciones de calentamiento y 300
muestras por variante (1\,200 observaciones; CSV de 1\,201 líneas).

\begin{center}\footnotesize
\begin{tabular}{@{}lrrrr@{}}
\toprule
\textbf{Garantía} & \textbf{p50 (ms)} & \textbf{p95 (ms)} & \textbf{p99 (ms)} & \textbf{ops/s} \\
\midrule
\texttt{w1} & 1,13 & 1,67 & 2,07 & 859 \\
\texttt{w1\_journal} & 2,26 & 3,89 & 5,70 & 402 \\
\texttt{majority} & 3,67 & 4,97 & 6,37 & 273 \\
\texttt{w3} & 2,57 & 3,10 & 3,39 & 382 \\
\bottomrule
\end{tabular}
\end{center}

La predicción previa situaba \texttt{w1} como la más rápida y
\texttt{majority} en torno a 1,5--3 veces su latencia. La mediana de
\texttt{majority} resultó 3,25 veces la de \texttt{w1}; exigir journal en
el primario elevó la mediana a 2,26 ms. Aunque \texttt{w3} fue más rápido
que \texttt{majority} en esta corrida, no representa una ventaja general:
son variantes medidas sucesivamente, con efectos de carga, caché y
sincronización. En tres nodos, \texttt{w3} exige tres confirmaciones y
\texttt{majority} dos; con latencia entre regiones la espera de réplicas
puede crecer mucho más. Mediciones realizadas en la máquina
descrita arriba, bajo contenedores locales (sin latencia de red interregional);
la memoria disponible para el motor Docker fue de 7,60 GiB, según docker info.

\textbf{Decisiones de los cinco TODO.} En \textbf{TODO 1},
\texttt{ordered=False} permite continuar el lote tras errores individuales;
en \textbf{TODO 2}, el índice compuesto ordena igualdad antes de rango;
\textbf{TODO 3} usa descuento condicionado y atómico;
\textbf{TODO 4} compara latencias por garantía; \textbf{TODO 5}, clasifica escrituras confirmadas, inciertas y pérdidas para
estimar RPO, sin confundir error del cliente con ausencia de escritura.

\section{Fallas (pasos 10 a 13)}

\textbf{Máquina de los pasos 10 a 13.} Intel Core i7-6500U (4 hilos), 11 GiB de RAM, Ubuntu 24.04.5
(kernel 7.0), Docker 29.9.2 sin límites de CPU ni RAM y \texttt{-{}-wiredTigerCacheSizeGB 0.25} por nodo.
\textbf{Se usó MongoDB 7.0.45 y no 8.0:} \texttt{mongo:8.0} no arranca en kernels 6.19 o más nuevos
(SERVER-121912, TCMalloc y \texttt{rseq}), así que un \texttt{docker-compose.override.yml} local (no
versionado) cambió solo la imagen. Replicación, elecciones, \texttt{w:majority},
\texttt{electionTimeoutMillis} y el \emph{step down} funcionan igual en 7.0.

\textbf{TODO 5: cómo se clasifica el RPO.} \texttt{calcular\_rpo} solo usa operaciones entre conjuntos.
\texttt{presentes} se lee con \texttt{readConcern: "majority"}: solo cuenta lo que tiene una mayoría y,
por tanto, no se puede revertir en un \emph{rollback}. Autotest: 3/3 \texttt{OK}; línea base
(\texttt{rpo-baseline.txt}): 20 intentos, 20 confirmados, 20 presentes, \texttt{PERDIDOS 0}.

\begin{center}\footnotesize
\begin{tabular}{@{}lll@{}}
\toprule
\textbf{Clave} & \textbf{Fórmula} & \textbf{Significado} \\
\midrule
\texttt{perdidos} & \texttt{confirmados} $-$ \texttt{presentes} & se confirmaron y no están: RPO $>$ 0 \\
\texttt{desconocidos\_presentes} & \texttt{sin\_confirmar} $\cap$ \texttt{presentes} & dieron error pero \textbf{sí} quedaron escritas \\
\texttt{desconocidos\_ausentes} & \texttt{sin\_confirmar} $-$ \texttt{presentes} & dieron error y no quedaron \\
\texttt{inesperados} & \texttt{presentes} $-$ \texttt{confirmados} $-$ \texttt{sin\_confirmar} & nunca se intentaron (error de la prueba) \\
\bottomrule
\end{tabular}
\end{center}

\begin{center}\footnotesize
\begin{tabular}{@{}llrrrlrr@{}}
\toprule
\textbf{Escenario} & \textbf{Caídos} & \textbf{Intentos} & \textbf{ok} & \textbf{error} & \textbf{Tipos de error} & \textbf{RTO} & \textbf{RPO} \\
\midrule
A, secundario & mongo-2 & 80 & 80 & 0 & --- & sin corte & 0 \\
B, primario & mongo-1 & 104 & 101 & 3 & \texttt{ServerSelectionTimeout} $\times$3 & 10,01 s & 0 \\
B5, \texttt{electionTimeoutMillis=5000} & mongo-3 & 113 & 112 & 1 & \texttt{ServerSelectionTimeout} $\times$1 & 4,99 s & 0 \\
C, sin mayoría & mongo-1 y 2 & 109 & 95 & 14 & \texttt{WTimeout} $\times$3, \texttt{WriteConcern} $\times$1, & $\approx$34 s* & 0 \\
 & & & & & \texttt{ServerSelectionTimeout} $\times$10 & & \\
\bottomrule
\end{tabular}
\end{center}
{\footnotesize * En C el primario vuelve cuando se reinician los nodos caídos, no por una elección automática:
ese «RTO» depende del operador (\texttt{sleep 20 + 10} del guion), no de MongoDB.}

\subsection{Falla A: cae un secundario}

\textbf{Predicción:} se esperaba un pico de latencia. \falta{resto de la predicción de la Falla A (¿errores?,
¿cuántos nodos quedan para la mayoría?), escrita antes de ejecutar (Persona C)}

\textbf{Precondición} (\texttt{falla-a-before.txt}): mongo-1 PRIMARY, mongo-2 y mongo-3 SECONDARY.
\textbf{Falla:} \texttt{docker stop -{}-timeout 0 ti4601-mongo-2}, señal \texttt{1792817808.142715757}
(2026-10-23T22:56:48.14-06:00, \texttt{falla-a-stop.txt}). \textbf{Durante} (\texttt{falla-a-status.txt}):
mongo-1 PRIMARY, mongo-2 \texttt{not reachable/healthy}, mongo-3 SECONDARY. Filas de \texttt{falla-a.csv}
(phase, seq, status, latency\_ms, completed\_epoch):
\begin{Verbatim}
antes-falla,12,ok,3.9,1792817807.919221
despues-falla,13,ok,5.2,1792817808.420788   <- 278 ms después de la señal
despues-falla,14,ok,4.7,1792817808.920558
\end{Verbatim}
Ningún error y ningún cambio apreciable de latencia: p50 de 5,3 ms antes (n = 13) y 4,7 ms después (n = 67);
p99 de 14,6 y 8,4 ms. Con \texttt{w:majority} el primario necesita su propio \emph{ack} más el de
\textbf{un} secundario; con mongo-2 caído quedan mongo-1 y mongo-3, 2 de 3 votos: hay mayoría para mantener
al primario y para confirmar escrituras. El pico esperado no se dio porque la escritura en curso ya esperaba
el \emph{ack} de mongo-3. Al volver (\texttt{falla-a-recovered.txt}), mongo-2 entra como SECONDARY y retoma
la replicación desde su último \texttt{optime} leyendo el oplog del primario; no necesita \emph{initial
sync} porque estuvo fuera unos 15 s, muy por debajo de la ventana del oplog. \texttt{falla-a-rpo.txt}: 80/80
presentes, \textbf{PERDIDOS 0}.

\subsection{Falla B: cae el primario}

\textbf{Predicción:} \falta{¿cuántos segundos sin primario?, ¿de qué parámetro depende?, ¿qué verá la sonda?,
escrita antes de ejecutar (Persona C)}

\textbf{Precondición} (\texttt{falla-b-before.txt}): mongo-1 PRIMARY, \texttt{term=1},
\texttt{electionTimeoutMillis=10000}. \textbf{Falla:} \texttt{docker kill ti4601-mongo-1}, señal
\texttt{1792817872.249831348} (2026-10-23T22:57:52.27-06:00). \textbf{Durante}
(\texttt{falla-b-status.txt}, 15 s después): mongo-1 caído, mongo-2 SECONDARY, \textbf{mongo-3 PRIMARY,
\texttt{term=2}}. Filas de \texttt{falla-b.csv}:
\begin{Verbatim}
antes-falla,11,ok,6.7,1792817871.591525      <- última confirmada antes de la falla
antes-falla,12,error,3009.0,1792817875.095516,ServerSelectionTimeoutError
despues-falla,13,error,3001.2,1792817878.097063,ServerSelectionTimeoutError
despues-falla,14,error,3000.9,1792817881.098176,ServerSelectionTimeoutError
despues-falla,15,ok,1161.6,1792817882.260059 <- primera confirmada tras la señal
\end{Verbatim}
\textbf{RTO} = 1792817882.260059 $-$ 1792817872.249831 = \textbf{10,01 s}. El seq 12 figura como
\texttt{antes-falla} porque la sonda lo envió con el \texttt{kill} en curso: la señal se escribe cuando
\texttt{docker kill} devuelve, unos cientos de ms después.

\textbf{Por qué el RTO supera \texttt{electionTimeoutMillis}.} Los secundarios dejan de ver los latidos
del primario (cada 2 s) y esperan \texttt{electionTimeoutMillis} (10 s, con una pequeña variación
aleatoria) antes de convocar la elección; luego viene la votación, con un \texttt{term} nuevo; después el
driver descubre al nuevo primario (\texttt{heartbeatFrequencyMS=1000} en la sonda), y por último se
reintenta la operación pendiente: el seq 15 tardó 1\,161,6 ms porque se emitió durante la elección.
\textbf{Errores y \texttt{retryWrites}:} el driver reintenta \textbf{una sola vez} y solo si encuentra un
primario dentro de \texttt{serverSelectionTimeoutMS=3000}; sin primario durante unos 10 s, tres intentos
seguidos agotaron los 3 s y salieron como \texttt{ServerSelectionTimeoutError}. El reintento oculta fallas
cortas, no una ventana mayor que ese plazo. \textbf{\texttt{term} y el viejo primario:} el \texttt{term} pasó
de 1 a 2; al volver, mongo-1 ve un \texttt{term} mayor y entra como SECONDARY
(\texttt{falla-b-recovered.txt}): todos tienen \texttt{priority} 1, mongo-3 está sano y un nodo con
\texttt{term} viejo no gana elecciones ni acepta escrituras. \textbf{RPO} (\texttt{falla-b-rpo.txt}): 101
confirmadas, 101 presentes, \textbf{PERDIDOS 0}; los 3 errores no quedaron escritos. Con
\texttt{w:majority} cada \texttt{ok} ya estaba en 2 de 3 nodos, y lo tenía al menos uno de los que eligieron
al nuevo primario.

\textbf{B5, opcional: \texttt{electionTimeoutMillis = 5000}} (\texttt{falla-b5-*}). Precondición: mongo-3
PRIMARY, \texttt{term=3}; tras \texttt{docker kill ti4601-mongo-3} el nuevo primario es mongo-2, con
\texttt{term=4}; al final se restauró el valor a 10\,000 (\texttt{falla-b5-recovered.txt}).
\begin{Verbatim}
antes-falla,12,ok,4.7,1792818096.629883
despues-falla,13,error,3095.6,1792818100.221038,ServerSelectionTimeoutError
despues-falla,14,ok,1593.0,1792818101.814277   <- primera confirmada
\end{Verbatim}
Con la señal en \texttt{1792818096.827114}, \textbf{RTO = 4,99 s}, frente a 10,01 s con el valor por
defecto. El RTO escala casi 1:1 con \texttt{electionTimeoutMillis}, que es el término dominante; bajarlo
acorta el corte a cambio de más elecciones innecesarias si la red tiene picos de latencia. RPO 0, 1 error no
escrito.

\subsection{Falla C: dos secundarios caídos, primario vivo (sin mayoría)}

\textbf{Predicción:} \falta{¿habrá primario?, ¿qué pasa con cada insert \texttt{w:majority}?, comparación con la
Falla B del Lab 1, escrita antes de ejecutar (Persona C)}

\textbf{Precondición} (\texttt{falla-c-before.txt}): mongo-3 PRIMARY, mongo-1 y mongo-2 SECONDARY; sonda con
\texttt{-{}-w majority -{}-wtimeout-ms 2000}. \textbf{Falla:} \texttt{docker kill ti4601-mongo-1
ti4601-mongo-2}, señal \texttt{1792817968.942381} (2026-10-23T22:59:28.94-06:00). \textbf{Durante}
(\texttt{falla-c-status.txt}, 20 s después): mongo-1 y mongo-2 caídos y \textbf{mongo-3 SECONDARY}: no queda
primario. Filas de \texttt{falla-c.csv}:
\begin{Verbatim}
antes-falla,13,ok,6.6,1792817959.666162                       <- última confirmada
antes-falla,14,error,2377.4,1792817962.537133,WTimeoutError
antes-falla,17,error,4583.9,1792817972.672220,WriteConcernError
despues-falla,18,error,3002.7,1792817975.675176,ServerSelectionTimeoutError
   ... seq 19 a 27: ServerSelectionTimeoutError, unos 3000 ms cada uno ...
despues-falla,28,ok,213.2,1792818002.918506                   <- primera confirmada
\end{Verbatim}
Los seq 14 a 17 figuran como \texttt{antes-falla} porque \texttt{docker kill} de dos contenedores tardó
varios segundos en devolver: los nodos ya estaban caídos desde \textasciitilde\texttt{...960}, pero la señal
se escribió en \texttt{...968.94}. Por los tiempos y el tipo de error, esas filas pertenecen a la falla.

\textbf{Por qué el sobreviviente baja a SECONDARY.} Un primario debe ver a una mayoría de votantes (2 de 3).
Sin latidos de los otros dos, mongo-3 hace \emph{step down} pasado el \texttt{electionTimeoutMillis}, aunque
siga vivo. Eso evita el \emph{split-brain}: un primario aislado podría aceptar escrituras que otra mayoría
nunca vería. No es una elección, es el primario renunciando; y como un solo nodo no gana ninguna elección,
nadie más asume. \textbf{Dos fases de error:} (1) \texttt{WTimeoutError} y \texttt{WriteConcernError} (seq 14
a 17): mongo-3 todavía era primario, \textbf{aplicó} la escritura localmente y no consiguió el \emph{ack} de
un segundo nodo dentro de \texttt{wtimeout=2000 ms}; el seq 17 terminó como \texttt{WriteConcernError} cuando
el primario hizo \emph{step down} con la operación esperando. (2) \texttt{ServerSelectionTimeoutError} (seq
18 a 27): ya no hay primario y nada llega al servidor.

\textbf{Error $\neq$ no escrito} (\texttt{falla-c-rpo.txt}):
\begin{Verbatim}
intentos 109 · confirmados 95 · con error 14 · presentes 99
PERDIDOS 0
errores que SÍ se escribieron ........... 4  [14, 15, 16, 17]
errores que NO se escribieron ........... 10
\end{Verbatim}
Los 4 \texttt{WTimeoutError}/\texttt{WriteConcernError} \textbf{sí están} en la base, leídos con
\texttt{readConcern: majority}: cuando volvieron mongo-1 y mongo-2 replicaron esas entradas del oplog de
mongo-3 y quedaron mayoritarias. Los 10 \texttt{ServerSelectionTimeoutError} nunca llegaron a un servidor.
\textbf{Ninguna escritura se confirmó sin mayoría} y el RPO sobre lo confirmado es 0. \textbf{Efecto en un
cliente que compra:} si \texttt{comprar\_entradas} recibe un \texttt{WTimeoutError}, concluye «no se compró»
y reintenta con un \textbf{nuevo} \texttt{orden\_id}, el descuento puede quedar duplicado, porque el primero
sí se aplicó. Por eso el RPO se mide sobre lo confirmado (lo que dio error tiene resultado desconocido hasta
consultar la base) y el P2 exige operaciones \textbf{idempotentes}: el reintento usa el \textbf{mismo}
\texttt{orden\_id} (\texttt{\_id} único, \texttt{DuplicateKeyError}, se devuelve el estado existente).

\textbf{Comparación con la Falla B del Lab 1} (CockroachDB, Raft). En los dos, sin quórum no hay escrituras
confirmadas (CP en CAP). En Raft el \emph{leaseholder} aislado tampoco confirma y la petición queda colgada o
con timeout; en MongoDB el primario aplica localmente y devuelve \texttt{WTimeoutError}, lo que deja
escrituras «huérfanas» que luego se replican (o se revierten con \emph{rollback} si otro nodo se hubiera
elegido). Para el cliente, los dos casos son ambiguos. \textbf{Recuperación}
(\texttt{falla-c-recovered.txt}): al reiniciar mongo-1 y mongo-2 vuelve la mayoría, mongo-3 gana la elección
y queda 1 PRIMARY + 2 SECONDARY.

\subsection{Mayoría 2 de 3}

\begin{center}\footnotesize
\begin{tabular}{@{}lrccc@{}}
\toprule
\textbf{Escenario} & \textbf{Votos vivos} & \textbf{¿Mayoría ($\geq$ 2)?} & \textbf{¿Primario?} & \textbf{¿Escrituras \texttt{w:majority}?} \\
\midrule
Normal & 3 & sí & sí & sí \\
A (un secundario caído) & 2 & sí & sí, el mismo & sí, sin corte \\
B (primario caído) & 2 & sí & sí, tras la elección (\textasciitilde 10 s) & sí, después del RTO \\
C (dos caídos) & 1 & \textbf{no} & \textbf{no} (\emph{step down}) & \textbf{no} \\
\bottomrule
\end{tabular}
\end{center}

Con 3 nodos se tolera \textbf{una} falla. La continuidad de A y de B se debe a que quedan 2 de 3 votos, que
siguen siendo mayoría; en C queda 1 y no hay escrituras confirmadas. Un cuarto nodo no aumentaría la
tolerancia: la mayoría de 4 es 3 y se seguiría tolerando una sola falla. De ahí el número impar de votantes.

\end{document}
